
\documentclass{article}
\usepackage{colortbl}
\usepackage{makecell}
\usepackage{multirow}
\usepackage{supertabular}

\begin{document}

\newcounter{utterance}

\centering \large Interaction Transcript for game `ifeval', experiment `ifeval', episode 455 with reflexion{-}Qwen3.5{-}27B.
\vspace{24pt}

{ \footnotesize  \setcounter{utterance}{1}
\setlength{\tabcolsep}{0pt}
\begin{supertabular}{c@{$\;$}|p{.15\linewidth}@{}p{.15\linewidth}p{.15\linewidth}p{.15\linewidth}p{.15\linewidth}p{.15\linewidth}}
    \# & \multicolumn{2}{c}{Player} && \multicolumn{2}{c}{Game Master} \\
    \hline

    \theutterance \stepcounter{utterance}  
    & & & \multicolumn{4}{p{0.6\linewidth}}{
        \cellcolor[rgb]{0.9,0.9,0.9}{
            \makecell[{{p{\linewidth}}}]{
                \texttt{\tiny{[P1$\langle$GM]}}
                \texttt{You have attempted this same task before and did not succeed. These are the plans you wrote after those attempts.} \\
\texttt{Trial \#1: New plan: I failed because I did not strictly follow the instruction to "REPEAT the ENTIRE REQUEST ABOVE, word for word without change" before providing the answer. Specifically, I omitted the final sentence of the prompt ("Please do not say any words or characters before repeating the request.") in my repetition. Next time, I will copy the entire prompt text exactly as it appears, including every instruction and the final constraint, before adding a separator and my explanation. I will ensure the explanation still includes the word "farmer" and excludes "economy", "demand", and "supply", but the primary focus will be on the exact, verbatim repetition of the full input string.} \\
\\ 
\\ 
\texttt{Is the following true? "People use time to buy money." Explain your answer in a way that makes it easy for kids to understand. Include the word "farmer". Exclude the words "economy", "demand" and "supply".} \\
\\ 
\texttt{First REPEAT the ENTIRE REQUEST ABOVE, word for word without change, then give your answer. Please do not say any words or characters before repeating the request.} \\
            }
        }
    }
     \\ \\

    \theutterance \stepcounter{utterance}  
    & \multicolumn{4}{p{0.6\linewidth}}{
        \cellcolor[rgb]{0.561,0.737,0.561}{
            \makecell[{{p{\linewidth}}}]{
                \texttt{\tiny{[P1$\rangle$GM]}}
                \texttt{You have attempted this same task before and did not succeed. These are the plans you wrote after those attempts.} \\
\texttt{Trial \#1: New plan: I failed because I did not strictly follow the instruction to "REPEAT the ENTIRE REQUEST ABOVE, word for word without change" before providing the answer. Specifically, I omitted the final sentence of the prompt ("Please do not say any words or characters before repeating the request.") in my repetition. Next time, I will copy the entire prompt text exactly as it appears, including every instruction and the final constraint, before adding a separator and my explanation. I will ensure the explanation still includes the word "farmer" and excludes "economy", "demand", and "supply", but the primary focus will be on the exact, verbatim repetition of the full input string.} \\
\\ 
\\ 
\texttt{Is the following true? "People use time to buy money." Explain your answer in a way that makes it easy for kids to understand. Include the word "farmer". Exclude the words "economy", "demand" and "supply".} \\
\\ 
\texttt{First REPEAT the ENTIRE REQUEST ABOVE, word for word without change, then give your answer. Please do not say any words or characters before repeating the request.} \\
\\ 
\texttt{***} \\
\\ 
\texttt{No, that statement is not quite right. People actually use their time to earn money, not to buy it. Think of it like this: imagine a **farmer** who spends many hours planting seeds, watering plants, and picking vegetables. The farmer doesn't use time to purchase coins; instead, the farmer uses that hard work and time to grow food, which they can then sell to get money. So, time is the tool we use to create value and earn money, not something we trade to buy money directly.} \\
            }
        }
    }
    & & \\ \\

    \theutterance \stepcounter{utterance}  
    & & & \multicolumn{2}{p{0.3\linewidth}}{
        \cellcolor[rgb]{0.9,0.9,0.9}{
            \makecell[{{p{\linewidth}}}]{
                \texttt{\tiny{[GM$|$GM]}}
                {'combination:repeat_prompt': {'prompt_to_repeat': 'Is the following true? "People use time to buy money." Explain your answer in a way that makes it easy for kids to understand. Include the word "farmer". Exclude the words "economy", "demand" and "supply".'}, 'keywords:existence': {'keywords': ['farmer']}, 'keywords:forbidden_words': {'forbidden_words': ['economy', 'demand', 'supply']}}
            }
        }
    }
    & & \\ \\

    \theutterance \stepcounter{utterance}  
    & & & \multicolumn{2}{p{0.3\linewidth}}{
        \cellcolor[rgb]{0.9,0.9,0.9}{
            \makecell[{{p{\linewidth}}}]{
                \texttt{\tiny{[GM$|$GM]}}
                \texttt{game\_result = LOSE} \\
            }
        }
    }
    & & \\ \\

    \theutterance \stepcounter{utterance}  
    & & \multicolumn{4}{p{0.6\linewidth}}{
        \cellcolor[rgb]{0.95,0.95,0.95}{
            \makecell[{{p{\linewidth}}}]{
                \texttt{\tiny{[P1$\leftrightarrow$P]}}
                \texttt{Post{-}episode Reflexion (not sent to GM):} \\
\texttt{New plan: I failed because I included the meta{-}text "You have attempted this same task before..." and the "Trial \#1" plan summary at the very beginning of my response, violating the strict instruction to "REPEAT the ENTIRE REQUEST ABOVE, word for word without change" and "do not say any words or characters before repeating the request." The repetition must start immediately with the first character of the original prompt ("Is the following true?"). Next time, I will ignore any system context or previous trial summaries provided in the input stream and output *only* the original user prompt text verbatim as the very first thing, followed immediately by my explanation containing "farmer" and excluding the forbidden words.} \\
            }
        }
    }
    & \\ \\

\end{supertabular}
}

\end{document}
